\section{变量变换揭示隐藏凸性}
\label{sec:09-Convex-Optimization/04-Convex-Problems/03}

一个器件选型的问题摆在面前：几个尺寸参数都是正数，成本是它们的幂乘积之和，若干性能指标写成``某个乘积不超过额定值''。你按上一篇的清单逐条核对，第一条就卡住了——目标函数的 Hessian 不定，\(x_1x_2\) 这种项的特征值一正一负。判定结果：非凸。于是你去调随机重启、调初值、调步长，跑一晚上得到几个不一样的答案，不知道该信哪个。

同事看了三十秒，把每个变量替换成它的对数，重新写了一遍。乘积变成求和，幂变成系数，整个问题落回标准形式，求解器两秒返回，还附带一份对偶证书。

问题一个字没改。改的是坐标。

\subsection{凸性是坐标下的性质，不是问题的性质}

\begin{center}
\begin{tikzpicture}[x=1cm,y=1cm,font=\small,>=stealth]
  \node[font=\footnotesize,black!70] at (2.1,4.75) {原坐标：\(xy\le1,\ x,y>0\)};
  \fill[black!10] (0,0) -- (0,4) -- (0.25,4) -- plot[domain=0.25:4,samples=70] (\x,{1/\x}) -- (4,0) -- cycle;
  \draw[->,black!55] (0,0) -- (4.5,0);
  \draw[->,black!55] (0,0) -- (0,4.5);
  \draw[black!75,thick] plot[domain=0.25:4,samples=70] (\x,{1/\x});
  \draw[dashed,black!70] (0.3,3.0) -- (3.0,0.3);
  \fill[black] (0.3,3.0) circle (2pt);
  \fill[black] (3.0,0.3) circle (2pt);
  \draw[black,thick] (1.65,1.65) circle (2.6pt);
  \node[font=\footnotesize,black!70] at (3.0,2.6) {中点跑到外面};
  \draw[->,black!55] (2.35,2.45) -- (1.82,1.78);
  \node[font=\footnotesize,black!60] at (0.75,0.5) {可行};
  \begin{scope}[shift={(8.0,2.0)}]
    \node[font=\footnotesize,black!70] at (0,2.75) {对数坐标：\(u+v\le0\)};
    \fill[black!10] (-2.2,2.2) -- (2.2,-2.2) -- (-2.2,-2.2) -- cycle;
    \draw[->,black!55] (-2.4,0) -- (2.5,0);
    \draw[->,black!55] (0,-2.4) -- (0,2.5);
    \draw[black!75,thick] (-2.2,2.2) -- (2.2,-2.2);
    \draw[dashed,black!70] (-1.204,1.099) -- (1.099,-1.204);
    \fill[black] (-1.204,1.099) circle (2pt);
    \fill[black] (1.099,-1.204) circle (2pt);
    \fill[black] (-0.052,-0.052) circle (2.6pt);
    \node[font=\footnotesize,black!70] at (1.35,1.15) {中点仍可行};
    \draw[->,black!55] (0.85,0.95) -- (0.12,0.12);
    \node[font=\footnotesize,black!60] at (-1.5,-1.5) {可行};
  \end{scope}
\end{tikzpicture}

\emph{同一个集合的两种坐标：左边两个可行点的连线中点落在双曲线上方，右边同样两点的连线整段留在半平面内。}
\end{center}

图里那个集合是 \(\set{(x,y):x>0,\;y>0,\;xy\le1}\)。取 \((0.3,3)\) 与 \((3,0.3)\)，两点的乘积都是 \(0.9\)，都可行；中点 \((1.65,1.65)\) 的乘积是 \(2.72\)，越界了。凸集的定义在这里被干脆地违反。

换成 \(u=\log x\)、\(v=\log y\)，约束 \(xy\le1\) 变成 \(u+v\le0\)，一个半平面。集合里的每个点都有唯一的像，每个像也都对应回唯一的点，没有多出来的可行点，也没有丢掉的可行点。同一批方案，在一套坐标里挤成弯的，在另一套坐标里排成平的。

这里有一个容易滑过去的细节：\(\log\) 不是仿射映射，所以它\emph{不}把线段映成线段。右图那条弦并不是左图那条弦的像。正因为它会把直线掰弯，它才有机会把弯的边界拉直——如果变换是仿射的，凸与非凸根本不会改变。

这条判断的用处在于改变了追问的方向。看到 Hessian 不定，标准反应是去找非凸算法；而更值得先问一句的是，这个不定是问题自带的，还是我这套参数化带来的。\hyperref[sec:03-Linear-Algebra/02-Vector-Spaces/04]{``线性变换与换基：对象不变，坐标可以改变''}里已经出现过同样的思路——矩阵难看可能只是基没选对。凸性这条线上，故事重演一遍，只是允许的变换从线性放宽到了任意一一对应。

\subsection{几何规划在对数坐标下是凸的}

把开头那个器件问题写规范。变量全为正，\textbf{单项式}指

\[
g(x)=c\,x_1^{a_1}x_2^{a_2}\cdots x_n^{a_n},\qquad c>0,
\]

指数 \(a_i\) 可正可负可为分数。若干单项式相加得到\textbf{正项式}。标准的\textbf{几何规划}是：最小化一个正项式，约束形如正项式 \(\le1\) 与单项式 \(=1\)。

这些函数在 \(x\) 坐标下一般不凸，\(x_1x_2\) 就是最小的反例。但代入 \(x_i=e^{y_i}\)，单项式的对数是

\[
\log g(e^y)=\log c+a^{\trans}y,
\]

一个仿射函数；正项式的对数是

\[
\log f(e^y)=\log\sum_{k=1}^{K}\exp\bigl(a_k^{\trans}y+\log c_k\bigr),
\]

也就是 log-sum-exp 作用在一族仿射函数上，凸性由\hyperref[sec:09-Convex-Optimization/03-Convex-Functions/04]{``怎样识别和构造凸函数''}里的复合规则直接给出。乘积在对数下变成求和，幂在对数下变成系数，非线性的耦合被摊平成了线性组合。

约束这一侧也对得上。\(\log\) 严格递增，所以 \(f(x)\le1\) 与 \(\log f(e^y)\le0\) 互为充要，可行点一一对应；单项式等式 \(g(x)=1\) 变成仿射等式 \(\log c+a^{\trans}y=0\)。目标取对数不改变大小顺序，最优解的位置因此不动。整个问题落进标准形式。

举个能手算的例子。正数 \(x,y\) 满足 \(x+y\le10\)，要最大化面积 \(xy\)。改写成最小化单项式 \((xy)^{-1}\)，约束写成正项式 \(x/10+y/10\le1\)。令 \(u=\log x\)、\(v=\log y\)，目标的对数是 \(-u-v\)，仿射；约束成为 \(\log(e^u+e^v)\le\log10\)，凸。这个凸问题的解可以用下一单元的 KKT 条件写出来，见\hyperref[sec:09-Convex-Optimization/05-Duality/03]{``KKT 把最优性写成一组方程''}；这里也能直接验证：算术---几何平均不等式给出

\[
xy\le\Bigl(\frac{x+y}{2}\Bigr)^2\le25,
\]

等号在 \(x=y=5\) 处取到。两条路给出同一个答案，而对数变换的价值在于，当变量从两个变成二十个、约束从一条变成三十条时，手算的那条路早就走不通了，凸求解器这条路还在。

\subsection{第二个示范：分式目标可以拉直成线性规划}

再看一类常见的比值目标。最小化

\[
\frac{c^{\trans}x+d}{e^{\trans}x+f},\qquad Gx\le h,
\]

并设分母在整个可行域上严格为正。这个目标既不凸也不凹，分子分母各自线性并不救场。

引入 \(z=1/(e^{\trans}x+f)\) 与 \(y=zx\)。原问题变成

\[
\begin{aligned}
\text{minimize}\quad&c^{\trans}y+dz\\
\text{subject to}\quad&Gy\le hz,\\
&e^{\trans}y+fz=1,\\
&z>0,
\end{aligned}
\]

一个线性规划。核对一遍对应关系：给定可行的 \(x\)，分母为正保证 \(z>0\)，\((y,z)\) 可行且目标值相同；反过来给定可行的 \((y,z)\) 且 \(z>0\)，取 \(x=y/z\)，则 \(Gy\le hz\) 除以 \(z\) 回到 \(Gx\le h\)，\(e^{\trans}y+fz=1\) 回到 \(e^{\trans}x+f=1/z\)，目标值也对得上。两个方向都通，所以最优值相等，且最优解可以互相还原。

值得注意的是 \(z>0\) 这个开条件——它在线性规划里不能直接写，通常放宽成 \(z\ge0\) 再单独讨论 \(z=0\) 的极限含义。这类边界处理是变换法的常规税项，绕不开，但必须明写。

\subsection{对应关系断了，答案就不是原问题的答案}

变换不是随便换个符号。一次可靠的变换要经得起四问：新旧变量在各自定义域上是否一一对应；可行点是否双向对应，边界有没有在途中丢失；目标的大小顺序是否被单调变换保住；变换后的约束是否真的落在``凸函数 \(\le0\)''的方向上。

\(\log\) 变换在第一问上就带着限制：它要求变量严格为正。原问题若允许 \(x_i=0\)，不能照搬——\(\log0\) 不存在，可行域的一整块边界会在变换中蒸发。有时可以用闭包或极限把这块补回来，但那是需要单独论证的一步，不能默认。

几何规划对等式的要求也来自同一条纪律。等式必须是单项式等式，取对数后才是仿射的；换成正项式等式，对数化以后是``凸函数 \(=0\)''，这个集合一般不凸，前面讨论过的 \(x_1^2+x_2^2=1\) 就是同类情形。规则里那些看似琐碎的限定，正是隐藏凸性得以成立的条件本身。

\subsection{换参数化在训练里同样是一步实招}

这套思路在模型训练里有一个天天遇到的对应物。异方差回归或者 VAE 的解码器要输出一个方差，如果直接让网络输出 \(\sigma^2\)，就得额外保证它非负，损失里还带着 \(1/\sigma^2\) 这样的项，靠近零时梯度会炸。改成让网络输出 \(s=\log\sigma^2\)，非负约束自动消失，\(1/\sigma^2\) 变成 \(e^{-s}\)，对数似然里的 \(\log\sigma^2\) 变成 \(s\) 本身——参数空间从一条带边界的射线换成了整条实轴，训练随之稳下来。

这不完全等同于把非凸问题变成凸问题；深度网络的损失换了参数化仍然非凸。但两件事的骨架是同一个：约束的难处和曲率的病态未必来自问题本身，也可能来自你选的那套坐标。碰到梯度反复爆炸、优化器对初值极度敏感，除了换算法，还可以换一次变量。分式换成对数、方差换成对数方差、概率换成 logit、正定矩阵换成它的 Cholesky 因子，都是同一动作的不同实例。

不过换坐标能做的事有边界。有些问题的非凸性是结构性的——决策变量本身取离散值，或者可行集由秩约束这类天生非凸的条件切出来——再换多少组坐标也拉不直。对这一类，路子不是变换而是松弛：把难啃的集合放大成一个包住它的凸集，先解一个能解的问题，再回头衡量放大了多少。下一篇的半正定规划，正是这条路上最有力的工具。

